\section{Architecture}
\label{sec:Architecture}

This chapter describes the overall system architecture, component layout and
major data flows in Archimesh.

\subsection{System purpose}

The system lets multiple Archi clients collaborate on the same model timeline
through a centralized server. It combines:

\begin{itemize}
\item a local Archi client plugin that captures and applies model changes
\item a Quarkus server that validates, orders, persists and broadcasts changes
\item Kafka for ordered model event fan-out
\item Neo4j for the materialized model graph, op-log, tags and merge metadata
\end{itemize}

The model timeline is intentionally linear: one \code{HEAD}, no branches,
immutable tags as named historical points.

Figure~\ref{fig:system-context} shows the client and server boundaries. Only
Quarkus connects to Neo4j and Kafka; Archi plugins and the browser use the
server's WebSocket and HTTP interfaces.

\manualfigure{system-context}{System components and transport boundaries.}{system-context}

\subsection{High-level flow}

\begin{enumerate}
\item A user edits a model in Archi.
\item The client plugin captures local EMF changes and converts them into
  Archimesh operations.
\item The server validates the op batch against preconditions, locks and
  idempotency state, then assigns authoritative revisions.
\item One Neo4j transaction commits the log, projection, head revision and
  durable publication intent. The submitting session receives acceptance.
\item The delivery worker publishes committed operations to Kafka in model
  revision order, retrying failures from durable state.
\item Other connected clients apply the accepted operations into their local
  Archi model.
\end{enumerate}

\subsection{Component layout}

\begin{longtable}{p{38mm}p{116mm}}
\toprule
Component & Responsibility \\
\midrule
\filepath{server/} & Quarkus server: WebSocket and REST endpoints, authorization,
  validation, revision assignment, idempotency, persistence, broadcast. \\
\filepath{client/org.gautelis.archimesh.plugin/} & Eclipse RCP plugin:
  WebSocket client, EMF change capture, operation mapping, remote apply,
  offline outbox, cache persistence. \\
\filepath{client/client-tests/} & Client plugin tests: CRDT merge, operation
  mapping, folder sync, outbox replay, notation inventory. \\
\filepath{admin-web/} & Svelte admin UI: model overview, catalog, versions,
  access control, sessions, audit. \\
\filepath{schemas/} & JSON schemas for common types, op batches, operations,
  lock events, presence and notation contracts. \\
\filepath{neo4j/} & Cypher schema for constraints and indexes, plus a
  documented graph model. \\
\filepath{scripts/} & Build, validation, development and operational scripts. \\
\filepath{docker-compose.yml} & Local dev infrastructure: Neo4j + Kafka. \\
\bottomrule
\end{longtable}

\subsection{Server component map}

The server is organized into distinct packages:

\begin{longtable}{p{42mm}p{112mm}}
\toprule
Package & Responsibility \\
\midrule
\code{endpoint} & WebSocket endpoint (\code{ArchimeshEndpoint}), REST
  endpoints (\code{ModelStateEndpoint}, \code{AdminEndpoint},
  \code{AdminUiEndpoint}, \code{ConsistencyEndpoint}). \\
\code{service} & Central coordination (\code{ArchimeshService}), plus
  service interfaces for validation, revision, locks, Neo4j repository,
  Kafka publish/consume, session registry and idempotency. \\
\code{CommitDeliveryService} & Writer-lease renewal and ordered retry of
  committed publication intents. \\
\code{CheckpointMaintenanceService} & Optional scheduled checkpoint
  creation, covered payload pruning and older-checkpoint retention. \\
\code{service.impl} & Concrete implementations: in-memory services for
  validation, revision, locks and sessions; Neo4j repository with split
  support classes for op-log, materialized state, reads, checkpoints and compaction;
  Kafka publisher and consumer implementations. \\
\code{auth} & Authorization: identity modes, role normalization,
  policy decision point (PDP), authorization service and transport. \\
\code{model} & Admin data classes: status, catalog, sessions, ACLs,
  export/import, compaction, integrity, audit. \\
\code{wire} & Wire protocol envelopes (\code{ClientEnvelope},
  \code{ServerEnvelope}) and message types for inbound
  (Join, SubmitOps, AcquireLock, ReleaseLock, Presence) and outbound
  (CheckoutSnapshot, CheckoutDelta, OpsAccepted, OpsBroadcast,
  LockEvent, PresenceBroadcast, Error). \\
\bottomrule
\end{longtable}

\subsection{Client plugin component map}

\begin{longtable}{p{42mm}p{112mm}}
\toprule
Package & Responsibility \\
\midrule
\code{ws} & WebSocket session manager, outbox handling, join/rejoin
  decisions, local submit tracking, inbound message dispatch. \\
\code{emf} & EMF change capture, operation mapping, remote op
  application, CRDT merge (LWW, OR-set, tombstones), model controller. \\
\code{startup} & Eclipse startup hooks, workbench lifecycle bridge,
  status line bridge. \\
\code{handlers} & Menu command handlers: connect, disconnect, open
  server model, switch active model, show last conflict, resynchronize. \\
\code{ui} & Connect dialog, server model catalog dialog, open server
  model dialog. \\
\code{notation} & Notation serializer and deserializer for Archi
  diagram fields. \\
\code{util} & Authentication hint helpers and lightweight JSON
  utilities. \\
\bottomrule
\end{longtable}

\subsection{Two kinds of state}

The system maintains two kinds of state at all times:

\begin{description}
\item[Local working state] Inside each Archi client. This is the EMF model
  the user sees and edits. It is not authoritative.
\item[Authoritative shared state] On the Archimesh server. The Neo4j graph
  is the source of truth. A verified checkpoint and complete operation tail,
  or the full log before the first checkpoint, provide recovery data.
\end{description}

The client plugin synchronizes between these two states: capturing local
changes, submitting them to the server, and applying remote changes into the
local model.

\subsection{Snapshot vs op-log}

The server keeps a materialized graph, accepted operation history and optional
recovery checkpoints in Neo4j. Their roles differ:

\begin{itemize}
\item The materialized graph makes joins, exports and admin inspection fast
  (no replay needed).
\item Retained batch payloads provide delta checkout and replay after a checkpoint.
\item A raw checkpoint preserves projection state and all merge metadata at
  a committed revision. Earlier covered payloads can then be pruned.
\item Compact acceptance receipts preserve batch identity and revision ranges
  after payload pruning. Tags retain independent historical snapshots.
\end{itemize}

Exports contain the current client-facing snapshot and either complete
history (v1), or a checkpoint, prefix receipts and complete tail (v2).
Checkpoints contain raw graph and merge state beyond the client snapshot.
Arbitrary historical replay before the recovery base requires an archived
full history; tagged snapshots remain available independently.

\subsection{Folder model}

Model-tree folders are first-class synchronized state with strict rules:

\begin{itemize}
\item Root folders have stable synthetic ids (e.g. \code{folder:root-business}).
\item User folders have immutable ids --- rename is allowed, but identity is
  stable.
\item Object placement in folders is explicit and persisted.
\item Cross-root folder moves are rejected by validation.
\item Non-empty folder delete is rejected.
\end{itemize}

This avoids drift between clients and preserves Archi model-tree structure.

\subsection{Linear timeline}

Every model has exactly one linear revision timeline with:

\begin{itemize}
\item One moving \code{HEAD}
\item Zero or more immutable tags
\item No branches
\end{itemize}

Clients submit batches against a \code{baseRevision}. The server assigns an
authoritative contiguous revision range. If a batch contains \(N\) ops it
consumes \(N\) revisions. Tags are read-only pointers to historical revisions,
stored with a snapshot of the model at that point.

Only \code{HEAD} is writable. Tagged revisions can be read but reject writes.

\subsection{Deployment and operational considerations}

\subsubsection{Authoritative writer and durability}

Run one authoritative server per Neo4j repository. A renewable database lease
rejects a second active writer; commit and maintenance transactions check the
lease, and timeline operations are serialized within the active service.
Active-active server operation is not supported. Locks, presence and live
session registries are process-local; committed state and publication intents
survive restart in Neo4j. See Section~\ref{sec:Messaging} for delivery and
Section~\ref{sec:Management} for checkpoint recovery.

\subsubsection{Transport security}

The WebSocket and REST endpoints do not terminate TLS themselves. In production,
deploy behind a reverse proxy (Nginx, Caddy, Traefik) that terminates TLS and
forwards to the Quarkus server. Example configurations are provided in
\filepath{server/examples/}.

When using \code{app.identity.mode=proxy}, the reverse proxy must also
forward trusted identity headers, and the WebSocket upgrade path must carry
those headers through the proxy to the server.

\subsubsection{Reconnection backpressure}

After a server restart, all previously connected clients may attempt to
reconnect simultaneously. To avoid overwhelming the server:

\begin{itemize}
\item Clients should apply jitter to reconnection attempts (randomized delay
  1--5 seconds before the first retry).
\item The server should reject joins that arrive faster than a configurable
  rate limit.
\item Admin diagnostics expose active session counts so operators can detect
  reconnection storms.
\end{itemize}

The current client implementation does not yet include reconnect jitter. This
is a known operational gap for production deployments with many concurrent
clients.

\subsubsection{Kafka single-partition requirement}

Each model's ops topic must have exactly one partition to preserve total
order. If multiple partitions are used, the total order guarantee breaks and
clients can receive operations in different orders, producing divergent state.
This is a hard constraint, not a recommendation.

If Kafka publication fails, the committed payload remains pending in Neo4j
and is retried by the delivery worker. The submitting client may already have
received acceptance while other clients await fan-out. There is no direct
broadcast bypass in the normal production delivery path.
